% REAL-DSL typesetting for the paper: one listing of R0 and inline macros that
% match it. Paste the PREAMBLE block into the preamble and the LISTING block
% where R0 is first mentioned (stage 1). Compiles with pdfLaTeX, IEEEtran.

% ---------------------------------------------------------------- PREAMBLE
\usepackage[T1]{fontenc}
\usepackage[scaled=0.86]{inconsolata}   % a clean typewriter face; 0.86 matches Times' x-height
\usepackage{upquote}                      % straight quotes " in \texttt and listings, never ``''
\usepackage{listings}
\usepackage{xcolor}
\definecolor{dslkw}{HTML}{1D1D1F}
\definecolor{dslstr}{HTML}{4A5161}
\definecolor{dslcmt}{HTML}{8A8F98}

% the DSL as a listings language: keywords exactly as the grammar spells them
\lstdefinelanguage{realdsl}{
  morekeywords={MAINTAIN,ACHIEVE,by,using,operationalized,as,performed,
    taking,input,producing,output,followed,if,in,scenario,where,assuming,ensuring},
  sensitive=true,
  morestring=[b]",
  morecomment=[l]{\#},
}
\lstset{
  language=realdsl,
  basicstyle=\ttfamily\scriptsize,
  keywordstyle=\bfseries\color{dslkw},
  stringstyle=\color{dslstr},
  commentstyle=\itshape\color{dslcmt},
  columns=fullflexible, keepspaces=true, upquote=true,
  breaklines=true, breakatwhitespace=true, breakindent=1.5em,
  xleftmargin=0.6em, aboveskip=3pt, belowskip=3pt, frame=none,
  numbers=left, numberstyle=\tiny\color{dslcmt}, numbersep=5pt,
  showstringspaces=false,
}

% inline: the same look as the listing, in running text
\newcommand{\dslkw}[1]{\textbf{\texttt{#1}}}          % keyword:   \dslkw{assuming}
\newcommand{\dsllit}[1]{{\color{dslstr}\texttt{#1}}}   % quoted string: \dsllit{"fog\_density <= 50"}
\newcommand{\dslid}[1]{\texttt{#1}}                    % identifier, no quotes: \dslid{fog\_density}

% ---------------------------------------------------------------- LISTING
% (abridged: the input/output declarations of the two tasks are omitted)
\begin{figure}[t]
\begin{lstlisting}[caption={The requirement $R_0$ of the study, in the REAL DSL
  (abridged: the tasks' input/output declarations are omitted). Lines 2--7
  name the system $S$, line 9 the scenario, lines 10--11 the domain
  assumptions $D$, lines 12--13 the soft goals in $R$.},label={lst:r0}]
MAINTAIN "Pedestrian Safety"
  by "Pedestrian Check" using "Perception Module"
     operationalized as
       "Detect Pedestrian" performed by "yolov5s"
  followed by "braking" using "braking module"
     operationalized as "Apply Brakes"
       if "pedestrian detection flag=True"
       performed by "proportional_braking"
  in scenario where "A pedestrian trying to cross the street."
  assuming "fog_density <= 50" & "sun_altitude_angle >= 0"
         & "initial_separation_m >= 15" & "pedestrian on foot"
  ensuring "resume_within_s <= 10"
         & "braking is smooth unless an emergency stop is needed"
\end{lstlisting}
\end{figure}

% ---------------------------------------------------------------- USAGE
% In the prose, refer to the listing instead of re-quoting the whole clause:
%   ... \dslkw{assuming} clauses (the domain $D$), e.g.\ \dsllit{"fog\_density <= 50"}
%   (Listing~\ref{lst:r0}, line~10), written over quantities the simulator can measure ...
% Operators stay ASCII inside DSL literals ("<=") because that is what the tool
% parses; use $\leq$ only in mathematics, never inside a DSL string.
